\chapter{바리온 생성, 우주의 매끄러움, 허블 마찰}

\lectureoverview{오늘 강의의 큰 흐름은 두 개의 작은 무차원 수로 정리된다. 첫째는 살아남은 물질--반물질 비대칭으로, 강의에서는 이를 \(10^8\)분의 1 정도로 어림한다. 이를 동역학적으로 설명하려면 바리온 수 위반, 입자--반입자 비대칭, 열평형에서의 이탈이라는 사하로프의 세 조건이 필요하다. 둘째는 분리 시점의 원시 밀도 대비로, 크기는 \(10^{-5}\) 정도다. 중력은 이런 요동을 지우기보다 증폭하므로, 요동이 작다는 사실은 인플레이션의 동기가 된다. 이어서 팽창 우주에서 균일한 스칼라장의 운동 방정식을 유도하고 \(3H\dot\phi\)가 허블 마찰임을 확인해 인플레이션의 역학을 준비한다.}

\section{엔트로피 문제를 뒤집어 묻기}

오늘 밤 인플레이션을 시작하겠지만, 먼저 바리온 생성을 마무리해야 한다. 바리온 생성이란 반물질보다 물질이 더 많이 남게 된 과정이다. 구체적인 수치 모형보다 그 이론적 틀이 더 확실하다. 그래서 핵심 수치가 제일원리에서 유도되지 않았는데도 이 주제를 간결하게 정리할 수 있다.

이 문제로 이어지는 역사적 계기는 1980년에 있었다. 밥 와고너가 사바스 디모풀로스와 레너드 서스킨드에게 우주의 엔트로피가 왜 그렇게 큰지 물었다. 우주론의 관점에서 바꾸어 말하면, 살아남은 양성자 하나당 광자가 왜 그렇게 많은지를 물은 것이다.

온도가 매우 높았을 때는 광자, 쿼크, 반쿼크, 전자, 양전자가 모두 풍부했고 열평형에 가까웠다. 쿼크와 반쿼크의 수밀도는 각각 광자 수밀도와 대략 같은 규모였다. 이 관찰을 보면 질문을 거꾸로 해 볼 수 있다. 소멸 뒤에 짝이 맞지 않는 물질이 왜 조금이라도 남았을까?

이 논증의 구조를 처음 제시한 공은 안드레이 사하로프에게 돌아간다. 디모풀로스와 서스킨드도 1980년에 이와 매우 가까운 조건들을 독립적으로 도출했고, 이후 역사적 기고문을 준비하다 사하로프의 1967년 논문을 발견했다. 이제 이 조건들은 마땅히 사하로프의 이름으로 불린다.

이제 엔트로피와의 관계를 명시할 수 있다. 흑체 복사의 엔트로피는 일정한 1차 규모의 계수를 제외하면 광자 수에 비례한다. 실제로
\begin{equation}
s_\gamma=\frac{4\rho_\gamma}{3T},
\qquad
n_\gamma\propto T^3,
\end{equation}
따라서 공변 부피 안의 열적 광자 기체에서는 \(S_\gamma\propto N_\gamma\)이다. 수치 계수 자체가 쟁점은 아니다. 문제는 광자 수와 양성자 수의 비가 엄청나게 크다는 점이다.

\begin{classroomqa}
\audiencequestion{광자의 양은 엔트로피와 어떤 관계가 있나요?}
\lecturerresponse{흑체 복사는 열복사이며, 그 엔트로피는 평균 광자 수에 비례합니다. 이 맥락에서 우주 엔트로피가 왜 큰지 묻는 것은 사실상 우주 마이크로파 배경에 양성자 하나당 광자가 왜 그렇게 많은지를 묻는 것입니다. 비례 관계에는 수치 계수들이 포함되지만, 이들은 차수 규모의 문제와는 무관합니다. \lecturetimestamp{00:04:49--00:05:55}}
\end{classroomqa}

우주가 식으면서 입자와 반입자는 소멸했고, 복사는 남아 있다가 결국 분리되었다. 그러므로 놀라운 점은 광자가 살아남았다는 사실이 아니다. 소멸 뒤에 짝이 맞지 않는 물질이 아주 조금 남았다는 점이다.

따라서 유용한 무차원 수는 다음과 같다.
\begin{equation}
\eta_B
\equiv
\frac{N_B-N_{\bar B}}{N_\gamma}
\sim
\frac{N_p-N_{\bar p}}{N_p+N_{\bar p}}
\sim 10^{-8},
\label{eq:c08-baryon-asymmetry}
\end{equation}
여기서 마지막 값은 정밀한 측정치가 아니라 강의에서 사용한 어림값이다. 특이해 보이는 수가 등장하면 ``원래 그런 값이었다''는 설명만으로는 만족스럽게 논의를 끝낼 수 없다. 정량적 메커니즘은 답이 왜 작고 왜 특정한 부호를 갖는지 설명해야 한다.

이 초과분을 \emph{물질}이라 하고 \emph{반물질}이라 하지 않는 데는 관례적인 면이 있다. 물질이란 우리가 이루어진 쪽이기 때문이다. 물리적으로 말하면 관측 가능한 영역 전체에서 한쪽 부호가 우세하다는 것뿐이다.

\section{바리온 수와 첫 번째 조건}

바리온 수는 양성자나 중성자가 \(B=1\)을 갖도록 정규화한다. 쿼크의 언어로 쓰면,
\begin{equation}
B=\frac{N_q-N_{\bar q}}{3}.
\label{eq:c08-baryon-number}
\end{equation}
인자 \(1/3\)는 중요하다. 보통의 바리온은 각각 원자가 쿼크 세 개를 포함하기 때문이다. 다른 강입자도 바리온 수를 가질 수 있지만, 불안정한 강입자는 결국 붕괴한다. 순 쿼크 초과량은 더 일반적인 장부 정리 수단이다.

\begin{classroomqa}
\classquestion{순 쿼크 수는 3을 곱해야 하나요, 나누어야 하나요?}
\lecturerresponse{3으로 나누어야 합니다. 양성자에는 쿼크가 세 개 있고 바리온 수가 1이므로 \(3B=N_q-N_{\bar q}\)입니다. 칠판의 수정은 정의의 일부이지, 물리적 논증을 바꾸는 것이 아닙니다. \lecturetimestamp{00:09:23--00:10:17}}
\end{classroomqa}

바리온 수는 보존될 수 있는 양이라는 점에서만 전하와 비슷하다. 전하는 장거리 쿨롱장을 만들지만 바리온 수는 그렇지 않다. 양성자는 전하를 띠므로 전기장을 만들지만, 중성자는 바리온 수를 가지면서도 순전하는 없다. 따라서 보존 여부와 힘을 만들어 내는 성질은 논리적으로 별개의 특성이다.

\begin{classroomqa}
\audiencequestion{바리온 수와 전하의 관련된 차이는 무엇인가요?}
\lecturerresponse{둘 다 장부 정리에 쓰이는 전하처럼 취급할 수 있지만, 장거리 전자기장을 만드는 것으로 알려진 것은 전하뿐입니다. 바리온 수는 이에 대응하는 쿨롱력을 만들지 않습니다. \lecturetimestamp{00:11:03--00:12:15}}
\end{classroomqa}

\(B\)가 정확히 보존된다면, 초기의 초과분은 영원히 고정된 채 남을 것이다. 입자 쌍의 생성과 소멸은 \(N_B\)과 \(N_{\bar B}\)를 각각 바꿀 수 있지만, 두 수를 같은 양만큼만 바꾸므로 그 차이는 변하지 않는다. 따라서 균형 잡힌 상태에서 시작하는 동역학 이론은 바리온 수가 변할 수 있어야만 작동할 수 있다.

가장 간단한 개략적 예는 양성자 붕괴다. 전하와 에너지 보존을 만족하는 다음 과정을 생각해 볼 수 있다.
\begin{equation}
p\longrightarrow e^+ +\gamma,
\qquad
\bar p\longrightarrow e^-+\gamma.
\label{eq:c08-decay-pair}
\end{equation}
이 반응들은 바리온 수 위반을 장부상으로 나타낸 모형이지, 실험에서 우세할 채널을 예측한 것이 아니다. 이 과정들은 전하를 보존하면서 \(B=+1\) 또는 \(-1\)를 0으로 바꾼다.

\begin{figure}[htbp]\centering
\begin{coursefitfigure}\begin{tikzpicture}[x=1cm,y=1cm,>=Latex]
\draw[->,thick] (0,1.7)--(2.0,1.7);
\draw[->,thick] (2.0,1.7)--(4.1,2.55);
\draw[->,thick] (2.0,1.7)--(4.1,0.85);
\node[above] at (0.9,1.7) {\(p,\ B=1\)};
\node[right] at (4.1,2.55) {\(\gamma\)};
\node[right] at (4.1,0.85) {\(e^+\)};
\draw[->,thick] (6.0,1.7)--(8.0,1.7);
\draw[->,thick] (8.0,1.7)--(10.1,2.55);
\draw[->,thick] (8.0,1.7)--(10.1,0.85);
\node[above] at (6.9,1.7) {\(\bar p,\ B=-1\)};
\node[right] at (10.1,2.55) {\(\gamma\)};
\node[right] at (10.1,0.85) {\(e^-\)};
\node at (2.0,0.25) {\(B:1\to0\)};
\node at (8.0,0.25) {\(B:-1\to0\)};
\end{tikzpicture}\end{coursefitfigure}\caption{전하 켤레인 개략적 반응 채널. 어느 과정이든 바리온 수를 위반하지만, 두 반응률이 같으면 어느 한쪽 부호도 선호되지 않는다.}\label{fig:c08-decay-redraw}\end{figure}

\begin{classroomqa}
\audiencequestion{양성자가 양전자와 중성미자로 붕괴할 수도 있나요?}
\lecturerresponse{페르미온 두 개로 이루어진 최종 상태는 초기 상태의 페르미온 홀짝성과 맞지 않습니다. 개략적인 최종 상태 \(e^++\gamma\)에는 페르미온이 하나 있으므로 대칭성 논증에 충분합니다. 현실적인 양성자 붕괴 이론에는 다른 입자와 채널이 들어갈 수 있지만, 여기서는 \(\Delta B\ne0\)를 보이는 것만이 목적입니다. \lecturetimestamp{00:14:39--00:15:07}}
\end{classroomqa}

\begin{classroomqa}
\audiencequestion{바리온 수의 정확한 보존은 대칭성에 해당하나요?}
\lecturerresponse{그렇습니다. 보존량은 연속 위상 대칭과 연관될 수 있습니다. 다만 그 대칭에 이름을 붙인다고 해서 자연이 그 대칭을 따른다는 사실이 별도로 증명되는 것은 아닙니다. 실제 판단 기준은 바리온 수를 바꾸는 과정이 일어나는지 여부입니다. 그러므로 보존 법칙만으로 대칭성을 추론한다면 논증은 순환적이 됩니다. \lecturetimestamp{00:15:08--00:15:46}}
\end{classroomqa}

\begin{classroomqa}
\audiencequestion{양성자--반양성자 소멸도 양성자 붕괴라고 볼 수 있나요?}
\lecturerresponse{용어는 부차적이지만, 보존 법칙은 그렇지 않습니다. \(p+\bar p\to\text{photons}\)에서는 초기 바리온 수가 \(1-1=0\)이고 최종 바리온 수도 0입니다. 소멸은 순 \(B\)를 바꾸지 않고 입자 두 개를 없애지만, 식~\eqref{eq:c08-decay-pair}의 한 입자 채널은 \(B\)를 바꿉니다. \lecturetimestamp{00:15:48--00:16:47}}
\end{classroomqa}

현재의 양성자는 믿기 어려울 만큼 안정적이다. 강의에서는 우주의 나이보다 훨씬 긴 \(10^{32}\)--\(10^{34}\)년 규모의 하한을 인용한다. 이 관찰은 정확한 보존을 입증하지 않는다. 현재 에너지에서 바리온 수를 바꾸는 과정이 극도로 느리다는 것만 말해 준다.

\section{이산 대칭과 비대칭의 필요성}

다음 질문은 법칙이 물질과 반물질을 구별하는가이다. 세 가지 이산 변환으로 답을 정리할 수 있다.
\begin{align}
C&:\ \text{particle}\longleftrightarrow\text{antiparticle},\\
P&:\ \mathbf x\longmapsto-\mathbf x,\\
T&:\ t\longmapsto-t.
\end{align}
여기서 \(P\)는 공간의 방향성을 뒤집고, \(T\)는 미시적인 시간 영상을 거꾸로 돌린다. 거시적인 열역학 제2법칙이 시간 반전 대칭처럼 보일 필요는 없다. 여기서 묻는 것은 바탕이 되는 미시적 동역학이다.

\begin{figure}[htbp]\centering
\includegraphics[width=0.82\textwidth]{lecture_08_figure_03.png}\caption{칠판에 수정해 적은 정의 \(B=(N_q-N_{\bar q})/3\)에서 전하 켤레 변환 \(C\)의 도입으로 이어진다. ‘전하’라는 단어 뒤에 미완성으로 남은 필기는 독립적인 증거로 사용하지 않는다.}\label{fig:c08-charge-board}\lectureframe{00:20:23.280}\end{figure}

상대론적 국소 양자장론은 \(C\), \(P\), \(T\)가 각각 대칭일 것을 요구하지 않는다. 하지만 표준적인 가정 아래에서는 이들을 결합한 변환 \(CPT\)를 요구한다. 어떤 과정이 가능하다면, 모든 입자를 반입자로 바꾸고 공간 배치를 반사한 뒤 운동 방향을 뒤집은 과정도 허용된다.

양자역학에서 시간 반전은 단순히 \(t\)의 부호를 바꾸는 것보다 더 미묘하다. 시간 반전 대칭인 퍼텐셜 안의 스핀 없는 파동함수에 대해,
\begin{equation}
i\hbar\frac{\partial\psi}{\partial t}=H\psi,
\qquad
T:\ \psi(\mathbf x,t)\longmapsto\psi^*(\mathbf x,-t).
\end{equation}
\(t\to-t\)를 바꾸면 시간 미분의 부호가 바뀐다. 복소켤레를 취하면 \(i\to-i\)도 바뀌어 방정식이 다시 성립한다. 스핀이 있는 입자의 경우 반유니터리 시간 반전 연산자는 더 복잡한 구조를 갖지만, 이 간단한 예에서도 복소켤레가 필요함을 확인할 수 있다.

\begin{classroomqa}
\audiencequestion{파동함수의 시간 반전에는 \(t\to-t\)뿐 아니라 복소켤레도 필요한가요?}
\lecturerresponse{그렇습니다. 슈뢰딩거 방정식에서 그 이유를 볼 수 있습니다. 시간 미분만 반전하면 방정식의 한쪽만 일관성 없이 바뀌지만, 복소켤레를 취하면 \(i\)의 부호도 함께 바뀝니다. \lecturetimestamp{00:23:58--00:25:04}}
\end{classroomqa}

이제 첫 번째 사하로프 조건을 말할 수 있다. 우주가 입자와 반입자가 같은 개수인 상태에서 시작했다면, \(B\)를 바꾸는 과정이 없는 한 입자 쌍의 생성과 소멸은 순 바리온 수 0을 보존했을 것이다.
\begin{equation}
\boxed{\text{baryon-number violation}.}
\end{equation}
이 조건이 없으면 대칭적인 초기 상태는 순 바리온 초과량을 얻을 수 없다.

현재 양성자가 안정적이라는 사실은 이 조건과 양립한다. 대통일 이론의 관점에서는 매우 무거운 입자, 질량 \(M_X\)인 입자가 바리온 수 위반을 매개할 수 있다. \(M_X\)보다 훨씬 낮은 에너지에서는 유효 진폭이 개략적으로 다음의 역수 거듭제곱에 비례해 억제된다.
\begin{equation}
\begin{aligned}
\mathcal A_{\Delta B\ne0}&\propto \frac{E^2}{M_X^2},\\
\Gamma_p&\sim \frac{m_p^5}{M_X^4}.
\end{aligned}
\label{eq:c08-heavy-suppression}
\end{equation}
붕괴율의 추정값은 결합 상수와 채널에 따른 계수를 제외하면 통상적인 차원 6 스케일링을 따른다. \(M_X\)가 매우 크면 차가운 양성자는 사실상 영구히 안정한 것처럼 보일 수 있지만, 충분히 뜨거운 초기 플라스마에서는 바리온 수를 바꾸는 반응이 가능하다.

\begin{classroomqa}
\audiencequestion{무거운 질량 \(M_X\)에 충돌 에너지도 이미 포함되어 있나요?}
\lecturerresponse{아닙니다. \(M_X\)는 매개 입자의 정지 질량 스케일이고, \(E\)는 충돌에서 이용할 수 있는 에너지입니다. \(E/M_X\)에 대한 저에너지 전개는 \(E\ll M_X\)일 때 강하게 억제되며, 열적 충돌 스케일이 무거운 입자의 문턱에 가까워질수록 억제가 약해집니다. 칠판에 휴리스틱하게 적은 분모는 유효장론의 관점에서 이해해야 합니다. \lecturetimestamp{00:32:32--00:34:02}}
\end{classroomqa}

초기 상태에 입자와 반입자가 같은 개수로 있었다고 하자. 바리온 수 위반만으로는 여전히 어느 한쪽도 선호되지 않는다. 붕괴 \(p\to e^++\gamma\)마다 전하 켤레인 짝반응 \(\bar p\to e^-+\gamma\)이 있고, 뜨거운 플라스마에서는 역반응도 일어난다. 미시적 확률이 같다면 바리온 수를 오가는 반응은 생기지만 평균적인 한쪽 방향의 흐름은 생기지 않는다.

\begin{classroomqa}
\audiencequestion{반양성자 반응도에서 바리온 수는 얼마나 변하나요?}
\lecturerresponse{들어오는 반양성자의 바리온 수는 \(B=-1\)이고, 전자와 광자의 바리온 수는 \(B=0\)이므로 이 채널에서 바리온 수는 \(-1\)에서 0으로 바뀝니다. 그 전하 켤레 반응에서는 \(+1\)에서 0으로 바뀝니다. 두 채널이 같은 빈도로 일어나면 어느 쪽도 순편향을 만들지 않습니다. \lecturetimestamp{00:37:24--00:37:39}}
\end{classroomqa}

무작위 요동은 관측된 초과량을 설명하기에는 훨씬 작다. 두 결과의 확률이 같은 독립 사건 \(N\)개가 있다면,
\begin{equation}
\Delta N\sim\sqrt N,
\qquad
\frac{\Delta N}{N}\sim N^{-1/2}.
\end{equation}
강의에서 어림한 관측 가능 우주 내 바리온 수 \(N\sim10^{80}\)를 대입하면
\begin{equation}
\frac{\Delta N}{N}\sim10^{-40},
\end{equation}
\(10^{-8}\)보다 32자릿수나 작다. 동전 던지기식 요동으로는 관측되는 모든 영역에서 같은 부호가 우세한 이유도 설명할 수 없다.

\begin{classroomqa}
\audiencequestion{관측 가능한 우주가 된 공변 영역에 처음부터 훨씬 많은 쌍, 이를테면 \(10^{160}\)쌍이 있었다면 \(10^{80}\)만큼 요동할 수 있지 않을까요?}
\lecturerresponse{제곱근 계산만 보면 그런 인위적인 수를 상정할 수는 있지만, 관측된 공간적 일관성은 설명하지 못합니다. 제한 없이 일어나는 통계적 요동이라면 보이는 우주 전체에 한쪽 부호만 나타나는 대신, 물질이 우세한 영역과 반물질이 우세한 영역이 자연스럽게 생길 것입니다. \lecturetimestamp{00:39:38--00:40:18}}
\end{classroomqa}

강의 당시의 관측 논증은 두 종류의 증거를 이용한다. 헬륨을 비롯한 보통의 원자핵은 우주선에서 관측되지만, 논의 대상이 된 자료에서는 반원자핵이 관측되지 않았다. 고에너지 충돌로 2차 반양성자가 생길 수 있으므로 반양성자만으로는 결정적이지 않다. 또한 물질 은하와 반물질 은하 사이의 경계에서는 뚜렷한 소멸 복사가 나와야 한다. 필요한 경계 신호가 없다는 점은 관측 가능 영역에 은하와 반은하가 같은 비율로 섞여 있다는 가설을 강하게 배척한다. 이는 강의의 역사적 맥락에서 제시한 관측 논증이며, 우리의 지평선 너머 영역에 관한 주장이 아니다.

\begin{classroomqa}
\audiencequestion{우주선에서 헬륨 원자핵도 실제로 측정하나요?}
\lecturerresponse{그렇습니다. 전하를 띤 헬륨 원자핵은 직접 측정됩니다. 우주선의 자세한 통계는 별도의 전문 분야지만, 흔한 원자핵과 그에 맞먹는 반원자핵 집단이 관측되지 않는다는 대조가 이 논증의 요점입니다. \lecturetimestamp{00:41:53--00:42:22}}
\end{classroomqa}

그러므로 미시적인 편향이 필요하다. 그것은 \(C\) 위반이며, 현실적인 바리온 생성 논의에서는 \(CP\) 위반도 필요하다. 중성 \(B\) 중간자 계는 실험실에서 볼 수 있는 사례를 제공한다. 중성 중간자와 그 반입자는 서로 전하 켤레인 최종 상태로 붕괴할 수 있는데, 부분 붕괴율은 서로 다를 수 있다.
\begin{equation}
\Gamma(B^0\to f)\ne\Gamma(\bar B^0\to\bar f).
\label{eq:c08-cp-partial}
\end{equation}
칠판에서는 정확한 채널과 비대칭 수치가 정해지지 않았다. 확실한 교훈은 약붕괴에서 입자--반입자 비대칭이 존재한다는 사실이 실험으로 확립되었다는 점이다.

그럼에도 \(CPT\) 정리는 진공에서 입자와 반입자의 질량과 총 붕괴 폭이 같아야 한다고 요구한다.
\begin{equation}
\Gamma_{\rm tot}(p)=\Gamma_{\rm tot}(\bar p).
\end{equation}
그렇다고 각 붕괴 분기비가 같아야 하는 것은 아니다. 한 쌍의 켤레 채널에서 나타나는 차이는 다른 채널, 이를테면 개략적인 전자 채널과 뮤온 채널의 차이로 상쇄될 수 있으므로 총합은 같게 유지된다.

\begin{classroomqa}
\audiencequestion{양성자와 반양성자의 총 붕괴율이 같다면 어떻게 편향이 생길 수 있나요?}
\lecturerresponse{총 붕괴 폭이 같다는 사실은 채널별 붕괴율까지 같다는 뜻이 아닙니다. 게다가 바리온 생성은 플라스마에서 일어납니다. 열적 점유 인자와 역반응, 산란 채널도 운동 방정식에 들어갑니다. 결정적인 점은 매질이 단순히 \(CPT\)를 회피한다는 것이 아니라, 상세 균형이 깨지면 \(CP\) 비대칭 부분 과정이 원천항을 만들 수 있다는 것입니다. \lecturetimestamp{00:49:39--00:51:48; 00:53:22--00:57:14}}
\end{classroomqa}

이것이 사하로프의 두 번째 조건이다.
\begin{equation}
\boxed{C\ \text{and}\ CP\ \text{violation}.}
\end{equation}
이 조건은 바리온 수를 바꾸는 반응에 방향성을 부여한다. 사하로프의 논문이 나오기 얼마 전 \(CP\) 위반이 발견되면서 이 조건은 물리적으로 구체화되었다.

최소 표준모형에 대해서는 세심하게 구분해야 한다. 표준모형은 보통의 섭동적 양성자 붕괴를 허용하지 않지만, 전약 이상 현상은 비섭동적으로 \(B+L\)를 위반할 수 있다. 고온에서는 스팔레론 전이가 그 위반을 효율적으로 일으킬 수 있다. 이것이 칠판에서 나중에 수정한 내용의 정확한 요지다. 표준모형의 원천만으로 우주 비대칭을 정량적으로 재현할 수 있는지는 별개의, 훨씬 어려운 문제다.

\section{열평형과 세 번째 조건}

장애물이 하나 더 남아 있다. 완전한 열평형에서는 모든 반응이 역반응과 균형을 이룬다. 바리온 수가 변할 수 있다면, 평형 상태의 화학 퍼텐셜은 임의의 초기 \(B\)를 보존할 수 없고, 상세 균형은 생성된 초과량을 지운다. 강의에서는 같은 요점을 \(CPT\)를 통해 설명한다. 정적인 평형 열욕에는 거시적인 시간의 화살이 없으므로 \(T\) 인자가 과정과 역과정의 차이를 만들어 줄 수 없다.

실제 우주는 팽창하며 식고 있다. 따라서 거시적 상태에는 시간의 순방향과 역방향을 구별하는 비대칭이 있다. 다만 팽창이 미시적 완화와 경쟁해 우세할 때만 이것이 중요하다.
\begin{equation}
\Gamma_{\rm reaction}\lesssim H.
\label{eq:c08-freezeout-condition}
\end{equation}
\(\Gamma_{\rm reaction}\gg H\)이면 충돌이 팽창에도 불구하고 평형을 빠르게 회복시킨다. 반응률이 허블률보다 작아지면 플라스마는 다시 조정될 수 없고, 풍부도는 동결된다.

\begin{classroomqa}
\audiencequestion{팽창만으로 계가 자동으로 평형에서 벗어나지 않는 이유는 무엇인가요?}
\lecturerresponse{아주 느리게 팽창하는 열욕은 미시 반응이 계속 상태를 조정하므로 임의로 평형에 가까운 상태를 유지합니다. 팽창 또는 냉각 시간 \(H^{-1}\)이 관련 반응 시간 \(\Gamma^{-1}\)보다 짧을 때만 평형에서의 이탈이 효과적으로 일어납니다. 그러면 거시적 상태를 시간에 따라 재생한 영상에서 어느 쪽이 순방향인지 드러납니다. \lecturetimestamp{01:01:22--01:02:26}}
\end{classroomqa}

이제 사하로프의 세 조건을 모아 쓸 수 있다.
\begin{enumerate}
\item 바리온 수 위반. 따라서 \(B\)가 변할 수 있다.
\item \(C\)와 \(CP\) 위반. 따라서 변화에 선호되는 부호가 생긴다.
\item 열평형에서의 이탈. 따라서 상세 균형이 그 선호를 상쇄할 수 없다.
\end{enumerate}
이 조건들은 여기서 논의하는 메커니즘에 필요하며, 초과량을 만드는 구조적으로 충분한 경로를 제공한다. 그러나 구체적인 고에너지 모형과 반응망, 우주론적 역사가 없이는 초과량의 크기를 정하지 못한다.

\begin{classroomqa}
\audiencequestion{측정된 비대칭 값이 고에너지 이론을 제약하나요?}
\lecturerresponse{원칙적으로는 그렇습니다. 실제로는 미지의 질량과 결합 상수, 위상, 반응률, 초기 우주의 여러 특성에 결과가 좌우됩니다. 따라서 상세한 모형을 정하지 않고 이 측정값 하나만으로 허용되는 범위를 좁히기는 어렵습니다. \lecturetimestamp{01:04:20--01:04:52}}
\end{classroomqa}

\begin{classroomqa}
\audiencequestion{바리온 비대칭과 전자 비대칭은 어떻게 연결되나요?}
\lecturerresponse{전기적 중성이 두 비대칭을 연결합니다. 개략적인 반응 채널에서는 양성자를 없애면 양전자가 생깁니다. 바리온 불균형을 반전시키면 렙톤 불균형도 이를 보상하는 방향으로 바뀝니다. 더 일반적으로는 전하 보존과 표준모형 반응을 통해 바리온 생성과 렙톤 생성이 서로 연결되지만, 보존되는 조합은 메커니즘에 따라 달라집니다. \lecturetimestamp{01:04:54--01:05:59}}
\end{classroomqa}

\begin{classroomqa}
\audiencequestion{관측된 표준모형의 \(CP\) 위반은 바리온 초과량을 설명하기에 너무 작은가요?}
\lecturerresponse{전약 바리온 생성에 관한 정밀 분석에서는 대체로 그렇게 결론 내립니다. 다만 강의에서는 계산이 모형에 의존하고 고에너지에서 미지수가 많다는 점을 강조합니다. 여기서 근거를 갖고 말할 수 있는 것은 저에너지에서 알려진 \(CP\) 위반이 존재하지만, 정량적으로 성공하는 바리온 생성 모형에는 대체로 추가 요소가 필요하다는 점입니다. \lecturetimestamp{01:06:01--01:07:50}}
\end{classroomqa}

\begin{classroomqa}
\audiencequestion{오늘날의 전자는 단지 반양성자가 소멸한 결과로 생겼나요?}
\lecturerresponse{개략적인 반응 채널만 보면 그런 장부 정리가 그럴듯해 보일 수 있지만, 초기 플라스마에는 렙톤도 별개의 거의 대칭적인 집단으로 존재했습니다. 관련 메커니즘은 바리온 생성이나 렙톤 생성으로 나타낼 수 있으며, 이후 반응들이 각 부문 사이에 비대칭을 전달합니다. 역사적으로 쓰인 \emph{바리온 생성}이라는 말을 모든 전자의 유일한 미시적 계보로 오해해서는 안 됩니다. \lecturetimestamp{01:07:51--01:08:44}}
\end{classroomqa}

\begin{classroomqa}
\audiencequestion{왜 모든 렙톤이 아니라 바리온을 중심으로 논의를 구성하나요?}
\lecturerresponse{바리온 물질은 원자 스펙트럼과 다른 천문 관측을 통해 조사할 수 있지만, 우주 중성미자 부문은 직접 세기가 훨씬 어렵습니다. 이 명칭에는 관측으로 접근할 수 있었던 대상과 이 분야의 역사도 반영되어 있습니다. \lecturetimestamp{01:08:45--01:09:12}}
\end{classroomqa}

\section{소멸 뒤에 남는 초과량}

고온에서는 입자 쌍 생성과 소멸이 양방향으로 일어난다.
\begin{equation}
\gamma+\gamma\longleftrightarrow p+\bar p,
\qquad
E_{\gamma\gamma}\gtrsim2m_pc^2\approx2\,\mathrm{GeV}.
\end{equation}
열욕이 식으면 바리온 쌍을 만들 만큼 에너지가 큰 광자는 지수적으로 드물어진다. 소멸은 여전히 에너지상 허용되지만, 역방향 생성은 더 이상 평형을 유지할 수 없다. 거의 모든 양성자가 반양성자를 만나고, 남는 것은 식~\eqref{eq:c08-baryon-asymmetry}의 작은 초과량이다. 소멸로 나온 에너지 높은 광자는 열욕과 빠르게 에너지를 주고받는다.

\begin{classroomqa}
\audiencequestion{비 \(10^{-8}\)은 현재의 값인가요? 큰 규모의 소멸은 언제 일어났나요?}
\lecturerresponse{이는 강의에서 제시한 현재 바리온 대 광자 비의 차수 규모입니다. 더 이른 시기에는 바리온과 반바리온이 모두 많았습니다. 역방향 쌍 생성이 효과를 잃자 대칭 성분이 소멸하고 순 초과량만 남았습니다. \lecturetimestamp{01:09:13--01:11:44}}
\end{classroomqa}

\begin{classroomqa}
\audiencequestion{양성자--반양성자 소멸은 언제 멈췄나요?}
\lecturerresponse{반양성자의 양이 무시할 수 있을 정도로 줄어들었을 때 멈췄습니다. 강의에서는 그 시점을 대략적으로만 특정합니다. 핵자합성보다 훨씬 이르고 재결합보다는 분명히 앞선, 기가전자볼트 규모의 온도에서 일어났습니다. 이 논의로부터 정확한 시각을 추론하지는 않습니다. \lecturetimestamp{01:11:45--01:12:13}}
\end{classroomqa}

소멸에서 나온 광자는 사라지지 않았다. 이 광자들은 결국 마이크로파 배경으로 이어지는 열복사의 역사에 속한다. 그런데도 오늘날 복사가 에너지 예산에서 차지하는 비중은 작다. 그 이유는
\begin{equation}
\rho_{\rm rad}\propto a^{-4},
\qquad
\rho_m\propto a^{-3}.
\end{equation}
광자 수는 많지만 광자 하나의 에너지는 적색편이로 줄어든다. 물질은 정지 질량 에너지를 유지한다.

\begin{classroomqa}
\audiencequestion{소멸로 광자가 그렇게 많이 생겼는데, 왜 오늘날 복사는 동역학적으로 중요하지 않나요?}
\lecturerresponse{광자는 주로 마이크로파 배경으로 여전히 존재하지만, 그 에너지 밀도에는 스케일 인자의 추가 거듭제곱이 반영됩니다. 팽창은 광자 수를 희석할 뿐 아니라 광자 하나하나의 에너지도 적색편이시킵니다. \lecturetimestamp{01:12:14--01:13:18}}
\end{classroomqa}

\begin{classroomqa}
\audiencequestion{매우 무거운 매개 입자가 있다면 표준모형이 불완전하다는 뜻인가요?}
\lecturerresponse{대통일 이론에서 양성자 붕괴를 매개하는 입자가 있다면, 이는 최소 표준모형을 넘어서는 새로운 물리일 것입니다. 표준모형의 이상에 의한 바리온 수 위반을 뒤에서 바로잡아 설명했다고 해서 그런 보통의 양성자 붕괴 입자가 생기는 것은 아니며, 그것만으로 정량적인 바리온 생성 모형이 완성되는 것도 아닙니다. \lecturetimestamp{01:13:20--01:13:31}}
\end{classroomqa}

\begin{classroomqa}
\audiencequestion{개략적인 두 붕괴 채널 사이에 큰 비대칭이 있으면, 처음에 \(p\)-\(\bar p\)만 있던 집단을 관측된 전자와 양성자의 같은 개수로 바꿀 수 있나요?}
\lecturerresponse{그런 장부 정리는 지어낼 수 있지만, 뜨거운 평형 플라스마에는 이미 접근 가능한 모든 입자와 반입자가 많이 들어 있습니다. 게다가 관측된 \(CP\) 비대칭이 한 채널을 임의로 우세하게 만들어도 된다는 뜻은 아닙니다. 일관된 모형이라면 칠판의 두 화살표만 볼 것이 아니라 결합된 전체 반응망의 변화를 계산해야 합니다. \lecturetimestamp{01:13:33--01:15:10}}
\end{classroomqa}

이 지점에서 바리온 생성에 관한 논의는 정직하게 끝난다. 정성적인 요소들은 알려져 있지만, 관측된 수치는 유도되지 않았다. 더 완전한 입자 이론과 초기 우주의 이력을 통제된 방식으로 기술해야 더 충분한 답을 얻을 수 있다.

\section{두 번째 수: 원시 밀도 요동}

인플레이션은 우주 마이크로파 배경을 통해 정량적인 문제가 된 또 다른 질문, 즉 큰 규모에서 우주가 왜 이렇게 균질하고 등방적인지를 다룬다. 우주 마이크로파 배경은 온도가 \(2.7\,\mathrm K\)에 가까운, 놀랄 만큼 정밀한 흑체 스펙트럼을 보이며 그 온도는 모든 방향에서 거의 같다. 한때 포괄적인 우주론적 인상에 머물렀던 것이 이제 고정밀 관측 사실이 되었다.

중력 때문에 이러한 균일성은 더욱 놀랍다. 확산은 잉크 한 방울을 퍼뜨려 농도를 균일하게 만든다. 인력인 중력은 정반대로 작용한다. 밀도가 주변보다 높은 영역은 주위 물질을 끌어당겨 밀도를 높이고 인접한 영역은 비워 낸다. 따라서 작은 섭동은 시간이 흐르면서 커져 은하와 은하단, 보이드를 만들어 낸다.

\begin{figure}[htbp]
\centering
\begin{coursefitfigure}
\begin{tikzpicture}[x=1cm,y=1cm,>=Latex]
\draw[thick] (-4.2,0) ellipse (1.25 and 0.72);
\draw[very thick] (0,0) ellipse (1.55 and 0.92);
\draw[thick] (4.2,0) ellipse (1.25 and 0.72);
\node at (-4.2,0) {underdense};
\node at (0,0) {overdense};
\node at (4.2,0) {underdense};
\draw[->,thick] (-2.65,0)--(-1.65,0);
\draw[->,thick] (2.65,0)--(1.65,0);
\draw[->,thick] (-2.85,0.75)--(-1.45,0.4);
\draw[->,thick] (2.85,0.75)--(1.45,0.4);
\draw[->,thick] (-2.85,-0.75)--(-1.45,-0.4);
\draw[->,thick] (2.85,-0.75)--(1.45,-0.4);
\end{tikzpicture}
\end{coursefitfigure}
\caption{중력 불안정성은 밀도가 낮은 주변 영역의 물질을 밀도가 높은 영역으로 옮겨 초기 대비를 확산시켜 없애는 대신 증폭한다.}
\label{fig:c08-gravity-growth}
\end{figure}

이 성장을 거꾸로 추적하면 초기 우주는 훨씬 더 균일했음을 알 수 있다. 그렇다고 완전히 균일해서도 안 된다. 대칭성이 정확히 유지된다면 은하는 형성될 수 없기 때문이다. 유용한 측도는 밀도 대비다.
\begin{equation}
\delta\equiv\frac{\delta\rho}{\rho},
\end{equation}
여기서 \(\delta\rho\)는 배경 밀도에서 국소적으로 벗어난 정도다. 짐 피블스와 초기 우주론자들이 제시한 구조 성장 추정에 따르면, 탈결합 당시 원시 진폭은 대략
\begin{equation}
\boxed{\frac{\delta\rho}{\rho}\sim10^{-5}}
\end{equation}
이었으며, \(10^{-4}\)--\(10^{-5}\)가 당시 인용된 범위다. 여기서도 이 값은 질문을 더 구체화할 뿐, 아직 그 값이 왜 그러한지는 설명하지 않는다.

앨런 구스가 1980에 제안한 인플레이션의 큰 틀은 엄청난 규모의, 거의 지수적인 팽창이 일어나는 시기다. 작고 매끄러운 영역이 여러 자릿수만큼 늘어나면 기존의 곡률은 희석되고 불균일성은 관측 가능한 영역보다 큰 규모로 밀려난다. 주름진 풍선의 비유는 매끄러워지는 직관을 잘 보여 주지만, 공간이 실제로 더 높은 차원에 놓여 있다고 받아들여서는 안 된다.

\begin{classroomqa}
\audiencequestion{주변 물질이 방향마다 눈에 띄게 들쭉날쭉한데, 우주는 어떻게 등방적일 수 있나요?}
\lecturerresponse{등방성은 큰 규모에서의 통계적 진술이지, 이 방이나 은하, 또는 모든 시선 방향이 똑같아 보인다는 뜻은 아니다. 지구는 국소적으로 산이 많지만 반지름에 비하면 매우 매끄럽다. 마찬가지로 은하와 은하단은 매우 매끄러운 큰 규모의 배경 위에 놓인 국소 구조다. \lecturetimestamp{01:25:37--01:26:22}}
\end{classroomqa}

\begin{classroomqa}
\audiencequestion{무작위 과정이라면 어쩌다 엄청나게 밀도가 높은 방향이 생길 수도 있지 않나요?}
\lecturerresponse{섭동은 제약 없이 서로 독립적으로 요동하는 것이 아니다. 거의 균질하고 등방적인 특정 통계적 스펙트럼을 가진 무작위 장을 이룬다. 파장이 길더라도 진폭은 작아서 \(10^{-5}\) 정도다. 스펙트럼과 그 기원은 앞으로 다룰 주제이지, 동전을 던질 때처럼 임의로 큰 편차가 생긴다고 가정하는 것이 아니다. \lecturetimestamp{01:26:26--01:28:23}}
\end{classroomqa}

\begin{classroomqa}
\audiencequestion{초기 우주의 높은 온도가 인플레이션을 일으켰나요?}
\lecturerresponse{인플레이션 시기에 온도는 주된 변수가 아니다. 핵심 현상은 장의 응력 에너지가 일으키는 급격한 지수 팽창이다. 인플레이션이 끝나고 그 에너지가 입자로 변환되면 열적 환경이 다시 중요해진다. \lecturetimestamp{01:28:25--01:29:04}}
\end{classroomqa}

\section{역학적 도입: 저항력이 작용하는 운동}

인플레이션 방정식을 준비하기 위해 점성 유체 속에서 낙하하는 단위 질량의 돌을 생각해 보자. 곧 돌을 대신할 스칼라장을 염두에 두고, 돌의 수직 좌표를 \(\phi\)라고 하자. 퍼텐셜 에너지가 \(V(\phi)\)라면 보존력은
\begin{equation}
F(\phi)=-\frac{dV}{d\phi},
\label{eq:c08-force-potential}
\end{equation}
이며, 저항력이 없을 때 뉴턴 방정식은 \(\ddot\phi=F\)이다.

\begin{figure}[htbp]
\centering
\includegraphics[width=0.82\textwidth]{lecture_08_figure_06.png}
\caption{낙하하는 돌의 비유를 나타내는 읽기 쉬운 판서 관계식 \(F=-dV/d\phi\). 아래쪽 화살표는 \(\phi\)가 증가하는 방향을 위쪽으로 정했음을 나타낸다.}
\label{fig:c08-force-board}
\lectureframe{01:35:09.513}
\end{figure}

선형 점성 저항력이 작용하면 추가되는 힘은 \(-\gamma\dot\phi\)이며, \(\gamma>0\)이다. 방정식은 다음과 같이 된다.
\begin{equation}
\ddot\phi+\gamma\dot\phi=F(\phi)
=-V'(\phi).
\label{eq:c08-viscous-motion}
\end{equation}
정지 상태에서는 저항력이 사라지고, 저항력은 언제나 속도의 반대 방향으로 작용한다. \(\phi\)가 위쪽을 향해 증가한다면 낙하하는 돌은 \(F<0\)이고 \(\dot\phi<0\)이다. 그러면 \(-\gamma\dot\phi>0\)이므로 저항력은 예상대로 위쪽을 향한다.

\begin{classroomqa}
\classquestion{퍼텐셜 힘과 저항력에 왜 같은 마이너스 부호가 붙나요? 두 힘은 서로 반대 방향이어야 하지 않나요?}
\lecturerresponse{두 힘의 방향은 서로 다른 양으로 정해진다. 퍼텐셜 힘은 \(-V'(\phi)\)이다. 저항력은 \(-\gamma\dot\phi\)이다. 아래쪽으로 움직이는 동안 \(F\)와 \(\dot\phi\)는 모두 음수이므로 \(F\)는 아래쪽을 향하고 \(-\gamma\dot\phi\)는 위쪽을 향한다. 따라서 두 식의 형태는 그대로 둔 채 대수적으로 서로 반대인 물리적 방향이 나온다. \lecturetimestamp{01:34:16--01:35:27}}
\end{classroomqa}

돌이 가속되면서 저항력은 두 힘이 균형을 이룰 때까지 커진다. \(F\)가 거의 일정할 때 종단 운동은 다음을 만족한다.
\begin{equation}
0=F-\gamma\dot\phi_{\rm term},
\qquad
\dot\phi_{\rm term}=\frac{F}{\gamma}.
\end{equation}
부호는 올바르다. 아래쪽을 향하는 \(F\)는 음의 종단 속도를 준다. 따라서 퍼텐셜이 완만하고 \(\gamma\)가 크면 돌은 경사면을 천천히 내려간다.

\begin{figure}[htbp]
\centering
\begin{tikzpicture}[x=1cm,y=1cm,>=Latex]
\draw[->] (0,0)--(7.0,0) node[right] {\(\phi\)};
\draw[->] (0,0)--(0,4.4) node[above] {\(V(\phi)\)};
\draw[very thick] plot[smooth] coordinates {(0.5,3.7) (1.6,3.55) (2.8,3.2) (4.0,2.55) (5.1,1.65) (6.3,0.6)};
\filldraw (1.45,3.58) circle (0.09);
\draw[->,thick] (1.55,3.48)--(2.25,3.23);
\node[align=left,anchor=west] at (2.45,3.75) {large drag,\\slow downhill motion};
\end{tikzpicture}
\caption{저항력이 큰 완만한 퍼텐셜. 장에 대응시키면 돌의 좌표는 공간적으로 균일한 스칼라 진폭으로 바뀐다.}
\label{fig:c08-shallow-potential}
\end{figure}

\section{인플라톤과 허블 마찰}

인플레이션 이론은 관례적으로 인플라톤이라 부르는 추가 스칼라장 \(\phi\)을 도입한다. 역사적으로 이는 우주의 매끄러움을 설명하기 위해 제안된 대담한 가설이었다. 공간적으로 거의 균일한 장 \(\phi=\phi(t)\)에서 시작하자. 팽창은 공간적 변화를 늘려 놓지만, 균일한 경우만 살펴보아도 핵심 동역학을 파악할 수 있다.

정준 스칼라장의 완전한 국소 에너지 밀도는
\begin{equation}
\rho_\phi
=
\frac12\dot\phi^{\,2}
+\frac{1}{2a^2}(\boldsymbol\nabla\phi)^2
+V(\phi).
\end{equation}
이다. 균일 근사 \(\boldsymbol\nabla\phi=0\)를 적용하면 다음과 같이 줄어든다.
\begin{equation}
\rho_\phi=\frac12\dot\phi^{\,2}+V(\phi).
\label{eq:c08-scalar-density}
\end{equation}
첫 번째 항은 장의 배치 공간에서의 운동에너지다. 이는 보통 공간을 통한 병진 운동이 아니라 시간에 따른 변화를 나타낸다.

\begin{classroomqa}
\audiencequestion{\(\phi\) 위의 점은 시간에 대한 편미분을 뜻하나요?}
\lecturerresponse{그렇다. 균일한 장에서는 \(\dot\phi=d\phi/dt\)이다. 공간 기울기가 있으면 별도의 양의 항이 더해지지만, 지금의 가정 아래서는 그 항이 사라진다. \lecturetimestamp{01:41:08--01:41:34}}
\end{classroomqa}

단위 공변 좌표 상자를 하나 따라가 보자. 그 물리적 부피는 \(a^3(t)\)이므로, 상자 안에 든 스칼라장의 에너지는
\begin{equation}
E_\phi
=a^3(t)
\left[
\frac12\dot\phi^{\,2}+V(\phi)
\right].
\end{equation}
이다. 명시적으로 들어간 인자 \(a^3(t)\)가 새롭게 추가된 요소다. 공변 좌표는 고정되어 있어도 상자의 물리적 부피는 팽창한다.

\begin{classroomqa}
\audiencequestion{척도인자 \(a\)는 공간 위치에 따라 달라지나요?}
\lecturerresponse{균일한 FRW 모형에서는 그렇지 않다. 같은 우주 시간에 공간 어디서나 값이 같지만 시간에 따라 변할 수 있다. \(\phi\)가 위치에 따라 달라진다면 공간 기울기에 따른 에너지를 다시 포함해야 한다. \lecturetimestamp{01:44:12--01:44:49}}
\end{classroomqa}

\begin{classroomqa}
\audiencequestion{이제 \(\phi(t)\) 하나만 쓰더라도 장의 값은 공간의 각 점과 연결되어 있나요?}
\lecturerresponse{그렇다. 장은 시공간의 모든 점에 값을 부여한다. 균일성이란 각 시각에 공간상의 값이 모두 같다는 뜻이며, 그 결과 동역학은 공통 함수 \(\phi(t)\) 하나로 표현된다. \lecturetimestamp{01:44:50--01:45:06}}
\end{classroomqa}

\begin{classroomqa}
\audiencequestion{장의 운동에너지 항은 어떤 물리적 운동을 나타내나요? 장은 더 강해지거나 약해지는 건가요?}
\lecturerresponse{그것은 운동 방정식과 초기값, 초기 속도, 그리고 \(V'(\phi)\)의 부호에 따라 정해진다. 이 비유는 에너지 공식만으로는 운동 방향을 알 수 없는 입자와 같다. \lecturetimestamp{01:45:07--01:46:08}}
\end{classroomqa}

균일한 장의 라그랑지안은 운동에너지에서 퍼텐셜 에너지를 뺀 값이다.
\begin{equation}
L_\phi=a^3(t)
\left[
\frac12\dot\phi^{\,2}-V(\phi)
\right].
\label{eq:c08-scalar-lagrangian}
\end{equation}
오일러--라그랑주 방정식을 적용하면
\begin{align}
\frac{d}{dt}\left(\frac{\partial L_\phi}{\partial\dot\phi}\right)
&=
\frac{\partial L_\phi}{\partial\phi},\\
\frac{d}{dt}\left(a^3\dot\phi\right)
&=-a^3V'(\phi).
\end{align}
곱의 미분법을 적용하면 두 항이 나온다.
\begin{equation}
a^3\ddot\phi+3a^2\dot a\,\dot\phi
=-a^3V'(\phi).
\end{equation}
\(a^3\)로 나누고 \(H=\dot a/a\)로 정의하면 다음을 얻는다.
\begin{equation}
\boxed{
\ddot\phi+3H\dot\phi+V'(\phi)=0
}
\label{eq:c08-hubble-friction}
\end{equation}
또는 \(F=-V'(\phi)\)를 사용하면
\begin{equation}
\ddot\phi+3H\dot\phi=F.
\end{equation}
각 단계는 식 ~\eqref{eq:c08-scalar-lagrangian}의 팽창 부피 인자에서 나온다.

\begin{classroomqa}
\audiencequestion{\(H=\dot a/a\)는 어디서나 같은 값인가요? 정말 상수인가요?}
\lecturerresponse{균일 모형에서는 같은 우주 시간 단면의 모든 공간점에서 값이 같다. 하지만 시간에 따라 일정할 필요는 없다. ``허블 상수''는 흔히 현재 값을 가리키며, \(H(t)\)는 진화 전반에 걸친 허블 팽창률을 뜻한다. \lecturetimestamp{01:50:32--01:50:57}}
\end{classroomqa}

식 ~\eqref{eq:c08-hubble-friction}은 이 비유를 정확하게 보여 준다. 점성 운동과 같은 형태이며, 저항 계수는
\begin{equation}
\gamma_{\rm cosmic}=3H.
\end{equation}
이다. 따라서 팽창은 균일한 스칼라장의 진화를 늦춘다. \(H\)가 크고 \(V(\phi)\)가 충분히 완만하다면 관성항은 저항항에 비해 작아지고, 장은 다음의 느린 구름 관계식에 가까워진다.
\begin{equation}
3H\dot\phi\simeq-V'(\phi).
\end{equation}
차갑고 점성이 큰 기름 속에서 공이 굴러 내려가듯, 장에 힘이 작용하더라도 완만한 퍼텐셜을 가로지르는 데 엄청나게 오랜 시간이 걸릴 수 있다.

다음 단계에서는 스칼라장의 에너지 밀도를 프리드만 방정식에 대입하여, 느리게 변하는 장에 팽창률이 어떻게 반응하는지 살펴본다. 이렇게 결합된 계가 인플레이션을 일으킨다. 지금까지 준비한 핵심은 다음과 같다. 장은 운동에너지와 퍼텐셜 에너지를 저장하고, 팽창하는 부피는 허블 마찰을 만들어 내며, 느리게 구르는 장은 급격한 팽창을 일으킬 만큼 오랫동안 퍼텐셜 에너지를 거의 일정하게 유지할 수 있다.

\section{기억해 둘 내용}

바리온 비대칭과 원시 밀도 대비는 서로 평행한 역할을 한다. \(10^{-8}\)와 \(10^{-5}\)는 이름만 붙인다고 설명되는 것이 아니지만, 각각 막연한 우주론적 인상을 정량적인 문제로 바꾼다. 전자는 사하로프 조건으로 이어지고, 후자는 인플레이션으로 이어진다.

인플레이션으로 이어지는 연결 고리도 명확하다. FRW 우주 속 균일한 스칼라장은
\[
\ddot\phi+3H\dot\phi+V'(\phi)=0.
\]
을 따른다. 가운데 항은 현상론적으로 임의로 추가한 점성이 아니다. 이 항은 물리적 부피 \(a^3(t)\)가 커지는 데서 나온다. 이 결과 하나가 다음 강의에서 퍼텐셜을 따라 천천히 구르는 장이 우주를 이끄는 과정을 설명하는 틀이 될 것이다.